\section{Optimal Control Problem}

Having established the well-posedness of the stochastic state equation, we now formulate the associated 
optimal control problem.

Let \(y_d\) denote a desired target state. The objective is to determine a control \(u\) that drives 
the system as close as possible to \(y_d\) while keeping the control effort moderate.

\subsection{Admissible Control Set}

We define the admissible control set by

\[
U_{ad}
=
L^2
\bigl(
\Omega;
L^2(0,T;H)
\bigr).
\]

More generally, box constraints may be imposed on the control, leading to the admissible set

\[
U_{ad}
=
\left\{
u\in L^2(\Omega;L^2(0,T;H))
;:;
u_a \le u \le u_b
\ \text{a.e. in } D\times(0,T)
\right\},
\]

where \(u_a\) and \(u_b\) are prescribed lower and upper bounds. In the present work, we 
restrict ourselves to the unconstrained case.

\subsection{Cost Functional}

The performance of a control \(u\) is measured through the quadratic cost functional

\begin{equation}
J(u)
=
\frac{1}{2}
\E
\left[
\int_0^T
|y(t)-y_d(t)|_H^2
,dt
\right]
+
\frac{\alpha}{2}
\E
\left[
\int_0^T
|u(t)|_H^2
,dt
\right],
\label{eq:cost}
\end{equation}

where

\begin{itemize}
\item \(y\) is the solution of the state equation corresponding to the control \(u\);
\item \(y_d\) is the desired state;
\item \(\alpha>0\) is a regularization parameter.
\end{itemize}

The first term in \eqref{eq:cost} measures the discrepancy between the actual state and 
the target state, while the second term penalizes excessive control effort.

\subsection{Optimization Problem}

The stochastic optimal control problem consists in finding

\[
u^\ast\in U_{ad}
\]

such that

\begin{equation}
J(u^\ast)
=
\min_{u\in U_{ad}}
J(u).
\label{eq:optimal_problem}
\end{equation}

The pair

\[
(y^\ast,u^\ast)
\]

satisfying \eqref{eq:optimal_problem} is called an optimal state-control pair.

\subsection{Existence of an Optimal Control}

We now establish the existence of at least one optimal control.

\begin{theorem}
Assume that

\[
y_d
\in
L^2
\bigl(
\Omega;
L^2(0,T;H)
\bigr).
\]

Then the optimization problem \eqref{eq:optimal_problem} admits at least one solution

\[
u^\ast\in U_{ad}.
\]
\label{thm:optimal_control}
\end{theorem}

\begin{proof}

Let

\[
{u_n}*{n\ge1}
\subset U*{ad}
\]

be a minimizing sequence such that

\[
J(u_n)
\longrightarrow
\inf_{u\in U_{ad}} J(u).
\]

Since the functional (J) contains the coercive term

\[
\frac{\alpha}{2}
\E
\int_0^T
|u_n(t)|_H^2,dt,
\]

the sequence \({u_n}\) is bounded in

\[
L^2
\bigl(
\Omega;
L^2(0,T;H)
\bigr).
\]

Consequently, there exists a subsequence, still denoted by \({u_n}\), and an element

\[
u^\ast
\in
U_{ad}
\]

such that

\[
u_n
\rightharpoonup
u^\ast
\]

weakly in

\[
L^2
\bigl(
\Omega;
L^2(0,T;H)
\bigr).
\]

Let \(y_n\) denote the corresponding sequence of state solutions. The a priori estimates 
established in Theorem~\ref{thm:existence} imply that \({y_n}\) is bounded in

\[
L^2
\bigl(
\Omega;
L^2(0,T;V)
\bigr).
\]

By weak compactness, there exists a subsequence converging to a limit \(y^\ast\), which 
is the state associated with \(u^\ast\).

Since the functional \(J\) is convex and weakly lower semicontinuous, we obtain

\[
J(u^\ast)
\le
\liminf_{n\to\infty}
J(u_n).
\]

Therefore,

\[
J(u^\ast)
=
\inf_{u\in U_{ad}}
J(u),
\]

which proves the existence of an optimal control.
\end{proof}

The existence result established above guarantees that the optimization problem is 
mathematically well posed. In the next section, we derive the adjoint equation and 
the first-order optimality conditions characterizing optimal solutions.
